\chapter{角动量耦合、外场效应与绝热演化}

\begin{mapbox}
本章补齐原子物理和量子材料面试中常被追问的三条链：角动量相加决定简并与选择定则；外场通过简并微扰解除简并；缓慢或周期变化的 Hamiltonian 决定绝热、共振和几何相位。默认使用 SI 制。
\end{mapbox}

\section{旋转、$SU(2)$ 与角动量相加}
绕单位矢量 $\vb n$ 转过无穷小角 $\dd\theta$ 的幺正算符写作
\[
U(\vb n,\dd\theta)=1-\frac{i}{\hbar}\dd\theta\,\vb n\cdot\vb J.
\]
连续旋转要求生成元满足
\[
[J_i,J_j]=i\hbar\epsilon_{ijk}J_k.
\]
定义 $J_\pm=J_x\pm iJ_y$，由对易关系逐项得到
\[
[J_z,J_\pm]=\pm\hbar J_\pm,\qquad
[J_+,J_-]=2\hbar J_z.
\]
对共同本征态 $\ket{jm}$，
\[
J_\pm\ket{jm}
=\hbar\sqrt{j(j+1)-m(m\pm1)}\ket{j,m\pm1}.
\]
升降必须终止，故 $m=-j,\ldots,j$，且 $2j$ 为非负整数。整数 $j$ 给 $SO(3)$ 的普通表示；半整数 $j$ 是其双覆盖 $SU(2)$ 的表示，转过 $2\pi$ 可多出负号。

两个独立角动量的总生成元为 $\vb J=\vb J_1+\vb J_2$。直积空间既可用
$\ket{j_1m_1}\ket{j_2m_2}$，也可用 $\ket{JM}$：
\[
\ket{JM}=\sum_{m_1m_2}
C^{JM}_{j_1m_1j_2m_2}\ket{j_1m_1}\ket{j_2m_2},
\qquad M=m_1+m_2.
\]
三角规则
\[
J=|j_1-j_2|,|j_1-j_2|+1,\ldots,j_1+j_2
\]
可由最高权态逐级剥离得到。

\subsection{$\tfrac12\otimes\tfrac12$ 的完整构造}
唯一的 $M=1$ 态为 $\ket{\uparrow\uparrow}$，所以
\[
\ket{1,1}=\ket{\uparrow\uparrow}.
\]
作用总降算符 $J_-=J_{1-}+J_{2-}$：
\[
J_-\ket{1,1}
=\hbar(\ket{\downarrow\uparrow}+\ket{\uparrow\downarrow})
=\hbar\sqrt2\ket{1,0}.
\]
因此
\[
\ket{1,0}=\frac1{\sqrt2}
(\ket{\uparrow\downarrow}+\ket{\downarrow\uparrow}),\qquad
\ket{1,-1}=\ket{\downarrow\downarrow}.
\]
$M=0$ 子空间剩余的归一化正交组合为
\[
\ket{0,0}=\frac1{\sqrt2}
(\ket{\uparrow\downarrow}-\ket{\downarrow\uparrow}).
\]
前三态交换对称，单态交换反对称。这一结论直接控制两个相同费米子的空间波函数对称性。

\subsection{$1\otimes\tfrac12$ 与一般 $l\otimes\tfrac12$}
最高权态 $\ket{\tfrac32,\tfrac32}=\ket{1,1}\ket{\uparrow}$。降一次并归一化：
\[
\ket{\tfrac32,\tfrac12}
=\sqrt{\frac23}\ket{1,0}\ket{\uparrow}
+\sqrt{\frac13}\ket{1,1}\ket{\downarrow}.
\]
与它正交的组合为
\[
\ket{\tfrac12,\tfrac12}
=\sqrt{\frac13}\ket{1,0}\ket{\uparrow}
-\sqrt{\frac23}\ket{1,1}\ket{\downarrow}.
\]
一般 $s=1/2$ 时
\[
\ket{j=l+\tfrac12,m}
=\sqrt{\frac{l+m+1/2}{2l+1}}\ket{m-\tfrac12,\uparrow}
+\sqrt{\frac{l-m+1/2}{2l+1}}\ket{m+\tfrac12,\downarrow},
\]
\[
\ket{j=l-\tfrac12,m}
=-\sqrt{\frac{l-m+1/2}{2l+1}}\ket{m-\tfrac12,\uparrow}
+\sqrt{\frac{l+m+1/2}{2l+1}}\ket{m+\tfrac12,\downarrow}.
\]
相位约定可以整体改变，但所有可观测量不变。基变换的幺正性给出 CG 系数的两组正交关系。

\section{Wigner--Eckart 定理与选择定则}
球张量 $T_q^{(k)}$ 满足
\[
[J_z,T_q^{(k)}]=\hbar qT_q^{(k)},\qquad
[J_\pm,T_q^{(k)}]
=\hbar\sqrt{k(k+1)-q(q\pm1)}T_{q\pm1}^{(k)}.
\]
它的矩阵元分解为
\[
\bra{\alpha'j'm'}T_q^{(k)}\ket{\alpha jm}
=\frac{\bra{\alpha'j'}T^{(k)}\ket{\alpha j}}
{\sqrt{2j'+1}}
C^{j'm'}_{jm\,kq}.
\]
约化矩阵元含径向和内部结构，CG 系数承担全部方向依赖。因此非零条件为
\[
m'=m+q,\qquad |j-k|\le j'\le j+k.
\]
电偶极算符 $\vb r$ 是 $k=1$ 的奇宇称张量，故
\[
\Delta l=\pm1,\qquad \Delta m=0,\pm1,
\]
且初末态宇称相反。选择定则只是“矩阵元必为零”的必要约束，并不保证允许跃迁强度一定大。

\section{从最小耦合到 Zeeman Hamiltonian}
电子电荷为 $-e$。均匀磁场取对称规范
$\vb A=\tfrac12\vb B\times\vb r$：
\[
H=\frac{1}{2m}(\vb p+e\vb A)^2+V(r)+g_s\mu_B\frac{\vb S}{\hbar}\cdot\vb B.
\]
因 $\nabla\cdot\vb A=0$，
\[
\vb p\cdot\vb A+\vb A\cdot\vb p=2\vb A\cdot\vb p
=(\vb B\times\vb r)\cdot\vb p=\vb B\cdot\vb L.
\]
故
\[
H=H_0+\frac{\mu_B}{\hbar}(\vb L+g_s\vb S)\cdot\vb B
+\frac{e^2}{8m}(\vb B\times\vb r)^2.
\]
弱场时最后的抗磁项是 $B^2$ 阶。若自旋可忽略，$\Delta E=\mu_BBm_l$，得到正常 Zeeman 效应。

\subsection{Landé 因子与弱场反常 Zeeman 效应}
自旋轨道耦合先把 $\vb L,\vb S$ 锁成 $\vb J$。定态中垂直于 $\vb J$ 的磁矩分量平均为零。利用
\[
\vb L\cdot\vb J=\frac12(J^2+L^2-S^2),\quad
\vb S\cdot\vb J=\frac12(J^2-L^2+S^2),
\]
把 $\vb L+g_s\vb S$ 投影到 $\vb J$，得到
\[
g_J=1+\frac{j(j+1)+s(s+1)-l(l+1)}{2j(j+1)}
\quad(g_s\simeq2),
\]
\[
\boxed{\Delta E=\mu_BBg_Jm_j}.
\]
当 $\mu_BB$ 大于自旋轨道能标，$\vb L,\vb S$ 分别绕磁场进动，$m_l,m_s$ 成为好量子数：
\[
\Delta E\simeq\mu_BB(m_l+2m_s).
\]
这就是 Paschen--Back 极限；在中间场区必须在固定 $m_j$ 子空间对角化完整 Hamiltonian。

\section{Stark 效应：宇称与简并的作用}
沿 $z$ 的静电场给电子的微扰
\[
H'=e\mathcal E z.
\]
若未扰态有确定宇称，则
\[
E_n^{(1)}=e\mathcal E\bra n z\ket n=0,
\]
因为被积函数为奇函数。非简并基态的最低修正因此为
\[
E_0^{(2)}
=e^2\mathcal E^2\sum_{m\ne0}
\frac{|\bra m z\ket0|^2}{E_0-E_m}
\equiv-\frac12\alpha\mathcal E^2.
\]
基态分母均为负，所以 $\alpha>0$，电场降低能量。

\subsection{氢原子 $n=2$ 线性 Stark 效应}
$n=2$ 的 $\ket{200}$ 与 $\ket{210}$ 等能，且 $z=r\cos\theta$ 只连接
$\Delta l=\pm1,\Delta m=0$。在这两个态构成的子空间中
\[
W=e\mathcal E
\begin{pmatrix}0&\bra{200}z\ket{210}\\
\bra{210}z\ket{200}&0\end{pmatrix}.
\]
分离径向和角向积分可得
\[
\bra{200}z\ket{210}=-3a_0
\]
（符号随相位约定改变）。所以本征态
\[
\ket\pm=\frac1{\sqrt2}(\ket{200}\pm\ket{210})
\]
的一阶能移为
\[
\boxed{\Delta E_\pm^{(1)}=\mp3ea_0\mathcal E}.
\]
$m=\pm1$ 的 $2p$ 态一阶不移动。线性 Stark 效应并未违反宇称论证：真正的零场本征态在简并子空间内并非必须选成宇称本征态。

\section{相互作用绘景、Dyson 展开与 Rabi 振荡}
令 $H=H_0+V(t)$，定义
\[
\ket{\psi_I(t)}=e^{iH_0t/\hbar}\ket{\psi_S(t)},\qquad
V_I=e^{iH_0t/\hbar}Ve^{-iH_0t/\hbar}.
\]
则
\[
i\hbar\pdv{}{t}\ket{\psi_I}=V_I(t)\ket{\psi_I}.
\]
积分后迭代：
\[
U_I(t)=1-\frac{i}{\hbar}\int_0^tV_I(t_1)\dd t_1
+\left(-\frac{i}{\hbar}\right)^2
\int_0^t\dd t_1\int_0^{t_1}\dd t_2\,V_I(t_1)V_I(t_2)+\cdots.
\]
时间有序的积分上限来自“较晚算符写在左边”，不能把二重积分随意换成方形区域而不加时间排序。

对能级差 $\hbar\omega_0$ 的二能级系统，近共振驱动在旋波近似下化为
\[
H_{\rm rot}=\frac{\hbar}{2}
\begin{pmatrix}-\Delta&\Omega\\ \Omega&\Delta\end{pmatrix},
\qquad \Delta=\omega-\omega_0.
\]
其本征频率为 $\Omega_R=\sqrt{\Omega^2+\Delta^2}$。初态为基态时
\[
P_e(t)=\frac{\Omega^2}{\Omega_R^2}
\sin^2\frac{\Omega_Rt}{2}.
\]
共振时振幅可达一；失谐既降低振幅又提高振荡频率。旋波近似要求驱动弱且
$|\Delta|,\Omega\ll\omega_0$。

\section{绝热定理与 Berry 相位}
设
\[
H(t)\ket{n(t)}=E_n(t)\ket{n(t)}
\]
且能级无简并。展开
\[
\ket{\psi(t)}
=\sum_n c_n(t)e^{-\frac{i}{\hbar}\int_0^tE_n(t')\dd t'}\ket{n(t)}.
\]
代入 Schrödinger 方程，动力学相位项抵消：
\[
\dot c_m=-\sum_n c_n\bra m\dot n\ket{}
e^{\frac{i}{\hbar}\int_0^t(E_m-E_n)\dd t'}.
\]
对本征方程求导并左乘 $\bra m$（$m\ne n$）：
\[
\bra m\dot H\ket n+(E_m-E_n)\bra m\dot n\ket{}=0,
\]
\[
\bra m\dot n\ket{}=\frac{\bra m\dot H\ket n}{E_n-E_m}.
\]
快速振荡可压低跃迁的充分判据是
\[
\boxed{
\left|\frac{\hbar\bra m\dot H\ket n}{(E_m-E_n)^2}\right|\ll1}.
\]
能隙闭合时该判据必然危险。保留 $m=n$ 项得到额外相位
\[
\gamma_n(t)=i\int_0^t\bra{n(t')}\dot n(t')\ket{}\dd t'.
\]
闭合参数回路 $C$ 上
\[
\gamma_n(C)=i\oint_C\bra n\nabla_{\vb R}n\ket{}\cdot\dd\vb R,
\]
这就是 Berry 相位；局域规范改变连接，但闭合回路相位模 $2\pi$ 不变。

\begin{interviewbox}
\textbf{30 秒回答绝热定理：}若 Hamiltonian 的变化时间尺度远长于能隙对应的量子时间，且演化中能隙不闭合，初态若是瞬时本征态，就保持在对应瞬时本征支上，只积累动力学相位和 Berry 几何相位。定量上非绝热耦合由
$\hbar\bra m\dot H\ket n/(E_m-E_n)^2$ 控制。
\end{interviewbox}

\begin{selftestbox}
\begin{enumerate}
\item 用总降算符构造 $\ket{\frac32,-\frac12}$。
\item 推导 $l\otimes\frac12$ 两组 CG 系数的归一化。
\item 为什么电偶极跃迁要求宇称改变？
\item 对 $l=1,s=1/2$ 分别计算 $j=3/2,1/2$ 的 $g_J$。
\item 为什么氢基态没有线性 Stark 效应，而 $n=2$ 有？
\item 从旋转框架 Hamiltonian 推导失谐 Rabi 公式。
\item 绝热演化经过真正能级交叉时会遇到什么问题？
\item Berry 连接为何规范依赖而 Berry 相位可观测？
\end{enumerate}
\end{selftestbox}
